
\documentclass{ieeeaccess}
\newlength{\xfigwd}
\usepackage{cite}
\usepackage{amsmath,amssymb,amsfonts}
\usepackage{algorithmic}
\usepackage{graphicx}
\usepackage{textcomp}
\usepackage{booktabs}
\usepackage{url}
\usepackage[T1]{fontenc}
\usepackage[scaled]{helvet}
\usepackage{adjustbox}
\renewcommand{\familydefault}{\sfdefault}

\def\BibTeX{{\rm B\kern-.05em{\sc i\kern-.025em b}\kern-.08em
    T\kern-.1667em\lower.7ex\hbox{E}\kern-.125emX}}
\begin{document}
\history{Date of publication xxxx 00, 0000, date of current version xxxx 00, 0000.}
\doi{10.1109/ACCESS.2017.DOI}

\title{SHIFT-ATCP: Shift-Aware and Diabetes-Protected Adaptive Conformal Prediction for Reliable Diabetes Classification Under Dataset Shift}
\author{\uppercase{Diwakar Diwakar}\authorrefmark{1}, \uppercase{Ansh Gupta}\authorrefmark{2}, AND
\uppercase{Gaurav Gupta}\authorrefmark{3}}

\address[1]{SRM University-AP (e-mail: diwakar.p@srmap.edu.in)}
\address[2,3]{Wenzhou University / Wenzhou Medical University / Wenzhou-Kean University (e-mail: ggupta@kean.edu)}

\tfootnote{This paragraph of the first footnote will contain support information, including sponsor and financial support acknowledgment.}

\markboth
{Diwakar \headeretal: SHIFT-ATCP for Reliable Diabetes Classification Under Dataset Shift}
{Diwakar \headeretal: SHIFT-ATCP for Reliable Diabetes Classification Under Dataset Shift}

\corresp{Corresponding author: Diwakar Diwakar (e-mail: diwakar.p@srmap.edu.in).}

\begin{abstract}
Machine learning models for diabetes prediction often perform well in the development dataset but become less reliable when transported to a new clinical population. This study evaluates external diabetes model transportability from the Kaggle Diabetes Prediction Dataset to NHANES 2017--2018 as the main external cohort and the Pima Indians Diabetes dataset as an extreme-shift stress-test cohort. Five model families were trained on the source cohort: logistic regression, random forest, XGBoost, LightGBM, and CatBoost. Three feature sets were studied: non-invasive screening features, routine clinical features, and extended transport features. To address reviewer concerns about calibration leakage and unclear split-specific metrics, all external datasets were split into target calibration and target test partitions using stratified 50/50 sampling with random seed 42. Calibration method selection used only source calibration or target calibration-select data, and final external metrics were computed only on held-out target test sets. The revised framework introduces SHIFT-ATCP, a shift-aware target-dominant adaptive conformal method that estimates source trust from Population Stability Index (PSI), domain-classifier AUC, target calibration size, and prevalence shift. It also introduces DP-SHIFT-ATCP, a diabetes-protected variant that selects class-specific conformal thresholds to improve diabetes-positive coverage. Dataset shift was severe: mean PSI was 1.545 for NHANES and 1.897 for Pima, and domain-classifier AUC was 1.000 for both targets. The shift-aware source trust mechanism therefore assigned near-zero source weights. Conventional unweighted ATCP undercovered both targets, with mean 90\% coverage of 0.8468 on NHANES and 0.7612 on Pima. Safety-gated SHIFT-ATCP restored near-nominal marginal coverage, reaching 0.9031 on NHANES and 0.9059 on Pima. DP-SHIFT-ATCP improved diabetes-positive coverage from 0.6612 to 0.9090 on NHANES and from 0.8092 to 0.9280 on Pima, with larger prediction sets and referral rates. These results show that external diabetes model reliability requires not only discrimination and calibration but also target-aware conformal uncertainty, source-trust estimation, and class-wise clinical safety analysis.
\end{abstract}

\begin{keywords}
diabetes classification, dataset shift, external validation, calibration, conformal prediction, SHIFT-ATCP, DP-SHIFT-ATCP, NHANES, Pima, clinical reliability.
\end{keywords}

\titlepgskip=-15pt
\maketitle

\section{Introduction}
Diabetes is a major public health problem and an important target for early screening and risk stratification \cite{who2023diabetes}. Machine learning models can combine demographic, clinical, and laboratory variables such as age, body mass index, hypertension, glucose, and HbA1c. However, high internal accuracy does not guarantee reliable use in another population. A model trained in one dataset may face different disease prevalence, different patient demographics, different measurement protocols, and different missingness patterns when transported to another dataset. This problem is known as dataset shift \cite{datasetShift}.

External transportability is especially important in clinical machine learning because wrong probability estimates or unreliable uncertainty estimates can affect clinical triage. A model may rank patients correctly but still produce poorly calibrated probabilities. Calibration methods such as Platt scaling, temperature scaling, isotonic regression, and beta calibration can improve probability reliability \cite{plattScaling,zadrozny2002,guo2017,kull2017}. However, calibration chosen only from the source population may not remain appropriate after transport.

Conformal prediction provides prediction sets with finite-sample marginal coverage under exchangeability \cite{vovk2005,shafer2008,angelopoulos2023}. In diabetes classification, the prediction set may be \{0\}, \{1\}, \{0,1\}, or empty depending on uncertainty. Under external dataset shift, the exchangeability assumption between source calibration and target test data is violated. This can cause conformal undercoverage, particularly for the minority diabetes-positive class.

This paper proposes a revised framework for reliable diabetes transportability. The source dataset is the Kaggle Diabetes Prediction Dataset. NHANES 2017--2018 is used as the main external validation cohort, and the Pima Indians Diabetes dataset is used as an extreme-shift stress test. The framework includes equal-budget model tuning, leakage-safe calibration selection, shift-aware conformal baselines, SHIFT-ATCP, safety-gated SHIFT-ATCP, and diabetes-protected DP-SHIFT-ATCP.

The main contributions are:
\begin{itemize}
    \item A leakage-safe external diabetes transportability protocol with fixed target calibration/test splits and final target-test-only reporting.
    \item A clear definition and computation of PSI using source-quantile bins, together with domain-classifier AUC as a second shift score.
    \item SHIFT-ATCP, a shift-aware conformal transport method that estimates source trust and down-weights source calibration when source-target shift is severe.
    \item DP-SHIFT-ATCP, a diabetes-protected conformal method that selects class-specific thresholds to improve diabetes-positive coverage.
    \item A clinical reliability analysis reporting marginal coverage, diabetes-positive coverage, set size, referral rate, clinical cost, bootstrap intervals, and target calibration-size ablation.
\end{itemize}

\section{Related Work}
External validation is necessary before clinical prediction models can be trusted in new populations \cite{tripod}. Dataset shift may occur as covariate shift, label shift, concept shift, or measurement shift \cite{datasetShift}. Diabetes datasets are particularly vulnerable because age, BMI, glucose measurement, HbA1c availability, diabetes prevalence, and outcome definitions vary across cohorts.

Probability calibration evaluates whether predicted risks match observed outcome frequencies. Common calibration approaches include Platt scaling, isotonic regression, temperature scaling, and beta calibration \cite{plattScaling,zadrozny2002,guo2017,kull2017}. Calibration is clinically relevant because a poorly calibrated probability can mislead threshold-based decisions even when discrimination is acceptable.

Conformal prediction gives prediction sets instead of only point predictions. Split conformal prediction uses a held-out calibration set to compute nonconformity thresholds \cite{vovk2005,shafer2008}. Under exchangeability, marginal coverage is guaranteed. However, external clinical deployment breaks this assumption when source calibration data and target test data come from different populations. Recent conformal transportability work therefore motivates target-aware, weighted, or adaptive conformal methods, especially when a small labeled target calibration set is available.

This work evaluates five tabular model families: logistic regression, random forest \cite{randomforest}, XGBoost \cite{xgboost}, LightGBM \cite{lightgbm}, and CatBoost \cite{catboost}. The focus is not simply on adding classifiers but on improving transport reliability, calibration, conformal coverage, and diabetes-positive clinical safety.

\section{Methods}
\subsection{Datasets and Feature Sets}
Three public datasets were used. The Kaggle Diabetes Prediction Dataset was used as the source development cohort \cite{mustafa2023dataset}. NHANES 2017--2018 was used as the main external validation cohort \cite{nhanes2018overview}. The Pima Indians Diabetes dataset was used as an extreme-shift stress-test cohort \cite{smith1988using}.

After cleaning and harmonization, the Kaggle source cohort contained 96,128 records with diabetes prevalence 8.82\%. NHANES contained 2,291 adults with diabetes prevalence 17.59\%. Pima contained 768 records with diabetes prevalence 34.90\%. Three feature sets were evaluated: Model A non-invasive screening features, Model B routine clinical features, and Model C extended transport features.

\begin{table*}[!t]
\centering
\caption{Dataset composition after cleaning and harmonization.}
\label{tab:dataset}
\begin{tabular}{lrrrrrr}
\toprule
Dataset & N & Class 0 & Class 1 & Diabetes \% & Mean age & Mean BMI \\
\midrule
Kaggle source & 96,128 & 87,646 & 8,482 & 8.82 & 41.80 & 27.32 \\
NHANES & 2,291 & 1,888 & 403 & 17.59 & 51.41 & 29.74 \\
Pima & 768 & 500 & 268 & 34.90 & 33.24 & 32.46 \\
\bottomrule
\end{tabular}
\end{table*}

\subsection{Split Construction and Leakage-Safe Evaluation}
The source dataset was split into source training, source calibration, and source internal test partitions. NHANES and Pima were each split into target calibration and target test partitions using a stratified 50/50 split with random seed 42. The target calibration and test splits were i.i.d. samples from the same cleaned target dataset, conditional on stratified sampling by diabetes status. The target test split was never used for model training, calibration method selection, source-weight selection, conformal threshold selection, or hyperparameter tuning.

\begin{table*}[!t]
\centering
\caption{Split construction and class distribution.}
\label{tab:splits}
\begin{tabular}{llrrrrl}
\toprule
Dataset & Split & N & Class 0 & Class 1 & Seed & Stratified by \\
\midrule
Kaggle & Source train & 57,676 & 52,587 & 5,089 & 42 & Diabetes \\
Kaggle & Source calibration & 19,226 & 17,529 & 1,697 & 42 & Diabetes \\
Kaggle & Source internal test & 19,226 & 17,530 & 1,696 & 42 & Diabetes \\
NHANES & Target calibration & 1,145 & 944 & 201 & 42 & Diabetes \\
NHANES & Target test & 1,146 & 944 & 202 & 42 & Diabetes \\
Pima & Target calibration & 384 & 250 & 134 & 42 & Diabetes \\
Pima & Target test & 384 & 250 & 134 & 42 & Diabetes \\
\bottomrule
\end{tabular}
\end{table*}

\subsection{Dataset Shift Metrics}
PSI was computed from numeric features and was not externally given. For each numeric variable, missing values were ignored. Ten bins were used. The bin edges were defined by source Kaggle quantiles. Source and target bin proportions were clipped at $10^{-6}$ before applying:
\begin{equation}
\mathrm{PSI}=\sum_b (p_{t,b}-p_{s,b})\log\left(\frac{p_{t,b}}{p_{s,b}}\right),
\end{equation}
where $p_{s,b}$ and $p_{t,b}$ are the source and target proportions in bin $b$. We also trained a domain classifier to distinguish source from target records. A domain AUC near 0.5 indicates similar source and target distributions, whereas a domain AUC near 1.0 indicates severe shift.

\begin{table*}[!t]
\centering
\caption{Dataset shift and shift-aware source trust.}
\label{tab:source_trust}
\begin{tabular}{llrrrr}
\toprule
Target & Feature set & Mean PSI & Domain AUC & Prevalence gap & Source weight \\
\midrule
NHANES & Screening & 1.227 & 1.000 & 0.087 & 0.000476 \\
NHANES & Routine clinical & 1.545 & 1.000 & 0.087 & 0.000252 \\
NHANES & Extended transport & 1.545 & 1.000 & 0.087 & 0.000252 \\
Pima & Screening & 1.861 & 1.000 & 0.261 & 0.000164 \\
Pima & Routine clinical & 1.897 & 1.000 & 0.261 & 0.000152 \\
Pima & Extended transport & 1.897 & 1.000 & 0.261 & 0.000152 \\
\bottomrule
\end{tabular}
\end{table*}

\subsection{Model Training and Hyperparameter Selection}
All models were trained only on the Kaggle source training partition. Five model families were evaluated for each feature set: logistic regression, random forest, XGBoost, LightGBM, and CatBoost. Equal-budget hyperparameter tuning was performed using source validation data only. Each model family used three candidate settings. The target datasets were not used for hyperparameter selection. The selected settings included $C=0.25$ or $C=4.0$ for logistic regression, depth and leaf constraints for random forest, depth and learning-rate choices for XGBoost and CatBoost, and leaf/learning-rate choices for LightGBM.

\begin{table*}[!t]
\centering
\caption{Selected equal-budget hyperparameter settings on source validation data.}
\label{tab:tuning}
\scriptsize
\begin{tabular}{lllccc}
\toprule
Feature set & Model & Selected setting & Val. ROC-AUC & Val. Brier & Val. ECE \\
\midrule
Screening & Logistic regression & $C=0.25$ & 0.832 & 0.176 & 0.276 \\
Screening & Random forest & depth=8, leaf=5 & 0.829 & 0.168 & 0.270 \\
Screening & XGBoost & depth=3, lr=0.05 & 0.834 & 0.068 & 0.003 \\
Screening & LightGBM & leaves=15, lr=0.05 & 0.830 & 0.171 & 0.258 \\
Screening & CatBoost & depth=4, lr=0.05 & 0.834 & 0.068 & 0.003 \\
Routine clinical & Logistic regression & $C=4.0$ & 0.964 & 0.080 & 0.132 \\
Routine clinical & Random forest & depth=10, leaf=5 & 0.976 & 0.060 & 0.129 \\
Routine clinical & XGBoost & depth=3, lr=0.05 & 0.980 & 0.022 & 0.003 \\
Routine clinical & LightGBM & leaves=15, lr=0.05 & 0.980 & 0.054 & 0.084 \\
Routine clinical & CatBoost & depth=6, lr=0.03 & 0.981 & 0.022 & 0.004 \\
Extended transport & Logistic regression & $C=4.0$ & 0.964 & 0.080 & 0.132 \\
Extended transport & Random forest & depth=10, leaf=5 & 0.977 & 0.060 & 0.128 \\
Extended transport & XGBoost & depth=3, lr=0.05 & 0.981 & 0.022 & 0.003 \\
Extended transport & LightGBM & leaves=15, lr=0.05 & 0.980 & 0.054 & 0.084 \\
Extended transport & CatBoost & depth=6, lr=0.03 & 0.981 & 0.022 & 0.004 \\
\bottomrule
\end{tabular}
\end{table*}

\subsection{Calibration and ECE Reporting}
ECE was computed using 10 equal-width bins over $[0,1]$:
\begin{equation}
\mathrm{ECE}=\sum_b \frac{n_b}{n}\left|\mathrm{acc}(b)-\mathrm{conf}(b)\right|.
\end{equation}
The revised protocol separates ECE by split. Source internal ECE is computed on the Kaggle internal test set. Source calibration method selection uses only the source calibration-select split. Adaptive target calibration selection uses only the target calibration-select split. Final external ECE is computed only on the held-out target test split. This avoids target-test leakage.

\subsection{SHIFT-ATCP and DP-SHIFT-ATCP}
For conformal prediction, the nonconformity score was defined as:
\begin{equation}
s_i=1-\hat{p}_{y_i}(x_i).
\end{equation}
Conventional ATCP combines source and target calibration scores. Under severe external shift, the large source calibration set can dominate the conformal threshold. SHIFT-ATCP therefore estimates a source-trust weight:
\begin{equation}
w_s=\exp(-2\overline{\mathrm{PSI}})\exp[-3\max(AUC_d-0.5,0)]n_t^{-1/2}\exp(-2\Delta_\pi),
\end{equation}
where $\overline{\mathrm{PSI}}$ is mean PSI, $AUC_d$ is domain-classifier AUC, $n_t$ is target calibration size, and $\Delta_\pi$ is absolute prevalence gap. Source and target scores are then combined using source weight $w_s$ and target weight 1.0. Safety-gated SHIFT-ATCP checks target calibration-select coverage and falls back to target-only conformal if validation coverage is below the nominal target.

DP-SHIFT-ATCP is a clinical safety variant. It selects separate class-specific conformal targets for non-diabetes and diabetes. The class-0 grid was 85\%, 87\%, 88\%, 89\%, and 90\%. The diabetes class-1 grid was 90\%, 92\%, 94\%, 95\%, and 97\%. The selected pair minimized a target calibration-select score combining marginal coverage error, diabetes undercoverage, class coverage gap, average set size, and clinical cost.

Uncertainty-efficiency frontiers were evaluated at confidence levels 80\%, 85\%, 90\%, 92\%, and 95\%. For each method, we reported marginal coverage, class-wise coverage, average prediction-set size, referral rate, diabetes prevalence among referred cases, non-referred error rate, and clinical cost. Bootstrap confidence intervals for coverage were computed by resampling the held-out target test predictions.

\section{Results}
\subsection{Dataset Shift and Source Trust}
Both external datasets showed severe shift. The average PSI was 1.545 for NHANES and 1.897 for Pima. The domain classifier AUC was 1.000 across target-feature-set combinations, showing that source and target records were highly distinguishable. As a result, the source-trust formula assigned near-zero source weights. This result is important: the method did not assume source calibration was useful under severe shift; it automatically suppressed source influence.

\begin{figure*}[!t]
\centering
\includegraphics[width=0.95\textwidth]{figures/figure_source_trust_stress_test.png}
\caption{Source-trust stress test under simulated source-target mixture shift. The top row shows that the estimated source-trust weight decreases as the target fraction increases, indicating reduced reliance on source calibration scores. The bottom row shows marginal and diabetes-positive coverage, with the dashed line indicating the nominal 90\% coverage target.}
\label{fig:stress}
\end{figure*}

\subsection{External Discrimination and Calibration}
On NHANES, the best external ROC-AUC was obtained by logistic regression with routine or extended clinical features and source-selected beta recalibration. ROC-AUC was 0.943, Brier score was 0.066, and ECE was 0.023. On Pima, the same logistic regression settings gave the best ROC-AUC of 0.806, Brier score 0.172, and ECE 0.040. These results show that simple calibrated logistic regression transported strongly for discrimination, but discrimination alone did not ensure conformal or class-wise reliability.

\begin{table*}[!t]
\centering
\caption{Best external target-test discrimination and calibration results.}
\label{tab:external_best}
\begin{tabular}{llllcccc}
\toprule
Target & Feature set & Model & Scenario & ROC-AUC & Brier & ECE & Sens./Spec. \\
\midrule
NHANES & Routine clinical & Logistic regression & Source-selected beta & 0.943 & 0.066 & 0.023 & 0.629/0.969 \\
NHANES & Extended transport & Logistic regression & Source-selected beta & 0.943 & 0.066 & 0.023 & 0.629/0.969 \\
Pima & Routine clinical & Logistic regression & Source-selected beta & 0.806 & 0.172 & 0.040 & 0.590/0.816 \\
Pima & Extended transport & Logistic regression & Source-selected beta & 0.806 & 0.172 & 0.040 & 0.590/0.816 \\
\bottomrule
\end{tabular}
\end{table*}

\subsection{Shift-Aware Conformal Reliability}
Conventional unweighted ATCP undercovered both targets. Mean 90\% coverage was 0.8468 on NHANES and 0.7612 on Pima. Weighted target-dominant ATCP improved coverage to 0.9023 on NHANES and 0.9017 on Pima. SHIFT-ATCP achieved 0.9029 on NHANES and 0.9059 on Pima. Safety-gated SHIFT-ATCP achieved 0.9031 on NHANES and 0.9059 on Pima. These results demonstrate that the main reason for earlier ATCP failure was excessive source influence under severe shift.

\begin{figure*}[!t]
\centering
\includegraphics[width=0.95\textwidth]{figures/figure_shift_aware_conformal_baselines.png}
\caption{Shift-aware conformal baseline comparison at 90\% confidence. The plot compares source-only conformal prediction, importance-weighted conformal prediction, target-only conformal prediction, conventional ATCP, weighted ATCP, SHIFT-ATCP, and clinically protected variants.}
\label{fig:shift_baselines}
\end{figure*}

\begin{table*}[!t]
\centering
\caption{Clinical reliability trade-off of SHIFT-ATCP variants at 90\% confidence.}
\label{tab:clinical_tradeoff}
\scriptsize
\begin{tabular}{llccccc}
\toprule
Target & Method & Marginal cov. & Diabetes cov. & Set size & Referral rate & Clinical cost \\
\midrule
NHANES & Old ATCP & 0.847 & 0.538 & 1.022 & 0.106 & 0.616 \\
NHANES & SHIFT-ATCP & 0.903 & 0.660 & 1.140 & 0.158 & 0.632 \\
NHANES & Safety-gated SHIFT & 0.903 & 0.661 & 1.140 & 0.158 & 0.632 \\
NHANES & DP-SHIFT & 0.885 & 0.909 & 1.220 & 0.220 & 0.619 \\
NHANES & Risk-asymmetric SHIFT & 0.898 & 0.940 & 1.286 & 0.286 & 0.716 \\
\midrule
Pima & Old ATCP & 0.761 & 0.484 & 1.197 & 0.249 & 1.393 \\
Pima & SHIFT-ATCP & 0.906 & 0.809 & 1.491 & 0.491 & 1.343 \\
Pima & Safety-gated SHIFT & 0.906 & 0.809 & 1.491 & 0.491 & 1.343 \\
Pima & DP-SHIFT & 0.875 & 0.928 & 1.475 & 0.475 & 1.176 \\
Pima & Risk-asymmetric SHIFT & 0.893 & 0.949 & 1.532 & 0.532 & 1.243 \\
\bottomrule
\end{tabular}
\end{table*}

\begin{figure*}[!t]
\centering
\includegraphics[width=0.95\textwidth]{figures/figure_clinical_reliability_tradeoff_shift_atcp_variants.png}
\caption{Clinical Reliability Trade-off of SHIFT-ATCP Variants. The figure separates marginal coverage, diabetes-positive coverage, average set size, and clinical cost for NHANES and Pima.}
\label{fig:clinical_tradeoff}
\end{figure*}

\begin{figure*}[!t]
\centering
\includegraphics[width=0.95\textwidth]{figures/figure_shift_atcp_coverage_efficiency_frontier.png}
\caption{Coverage-efficiency frontier for SHIFT-ATCP variants across confidence levels. The frontier shows the trade-off between higher coverage and larger prediction sets or referral burden.}
\label{fig:coverage_frontier}
\end{figure*}

\subsection{Diabetes-Protected Coverage}
Safety-gated SHIFT-ATCP gave strong marginal coverage, but diabetes-positive coverage remained lower than marginal coverage. DP-SHIFT-ATCP directly addressed this weakness. On NHANES, diabetes-positive coverage improved from 0.6612 for safety-gated SHIFT-ATCP to 0.9090 for DP-SHIFT-ATCP. On Pima, diabetes-positive coverage improved from 0.8092 to 0.9280. Risk-asymmetric SHIFT-ATCP provided even higher diabetes coverage, 0.9403 on NHANES and 0.9486 on Pima, but with larger set sizes and higher referral rates.

The best balanced individual DP-SHIFT result on NHANES used routine clinical random forest and achieved marginal coverage 0.9005, non-diabetes coverage 0.9015, diabetes coverage 0.8960, class coverage gap 0.0054, and average set size 1.1134. This shows that DP-SHIFT can identify operating points with balanced class coverage.

\begin{figure*}[!t]
\centering
\includegraphics[width=0.95\textwidth]{figures/figure_shift_atcp_clinical_referral_cost.png}
\caption{Clinical referral and cost profile for SHIFT-ATCP variants. Diabetes-protected methods increase referral and set size when needed, but reduce diabetes-positive undercoverage and clinical cost in the external target cohorts.}
\label{fig:clinical_referral_cost}
\end{figure*}

\subsection{Calibration-Size Ablation}
The calibration-size ablation tested NHANES calibration sizes 50, 100, 200, 500, 750, and 1145, and Pima calibration sizes 50, 100, 200, 300, and 384. Each setting used repeated stratified subsampling. At full calibration size, SHIFT-ATCP reached marginal coverage 0.9022 on NHANES and 0.9088 on Pima. DP-SHIFT reached diabetes-positive coverage 0.9221 on NHANES and 0.9419 on Pima. Therefore, standard SHIFT-ATCP is best interpreted as a marginal-coverage method, while DP-SHIFT is best interpreted as a diabetes-safety method.

\begin{table*}[!t]
\centering
\caption{Full target calibration-size ablation for selected conformal methods.}
\label{tab:ablation_full}
\scriptsize
\begin{tabular}{llccccc}
\toprule
Target & Method & Marginal cov. & Diabetes cov. & Class gap & Set size & Clinical cost \\
\midrule
NHANES & Old ATCP & 0.844 & 0.371 & 0.575 & 0.940 & 0.614 \\
NHANES & SHIFT-ATCP & 0.902 & 0.537 & 0.444 & 1.063 & 0.570 \\
NHANES & DP-SHIFT & 0.901 & 0.922 & 0.037 & 1.266 & 0.685 \\
NHANES & Risk-asymmetric SHIFT & 0.899 & 0.942 & 0.053 & 1.293 & 0.728 \\
\midrule
Pima & Old ATCP & 0.807 & 0.617 & 0.293 & 1.234 & 1.196 \\
Pima & SHIFT-ATCP & 0.909 & 0.796 & 0.173 & 1.480 & 1.336 \\
Pima & DP-SHIFT & 0.882 & 0.942 & 0.091 & 1.502 & 1.203 \\
Pima & Risk-asymmetric SHIFT & 0.897 & 0.951 & 0.083 & 1.543 & 1.257 \\
\bottomrule
\end{tabular}
\end{table*}

\begin{figure*}[!t]
\centering
\includegraphics[width=0.95\textwidth]{figures/figure_target_calibration_size_ablation_coverage.png}
\caption{Target calibration-size ablation for conformal coverage.}
\label{fig:ablation_coverage}
\end{figure*}

\begin{figure*}[!t]
\centering
\includegraphics[width=0.95\textwidth]{figures/figure_target_calibration_size_ablation_ece.png}
\caption{Target calibration-size ablation for expected calibration error. The plot shows how probability calibration stability changes as the number of labeled target calibration samples increases.}
\label{fig:ablation_ece}
\end{figure*}

\subsection{Source-Trust Stress Test}
A simulated source-trust stress test mixed source-like and target-like calibration/test data at target mixture levels of 0\%, 25\%, 50\%, 75\%, and 100\%. The source-trust weight decreased as target mixture and domain shift increased. For NHANES routine features, mean source weight decreased from 0.0283 at 0\% target mixture to 0.0003 at 100\% target mixture. For Pima routine features, it decreased from 0.0484 to 0.0002. This validates that the source-trust mechanism is not a fixed target-only rule; it responds to measured shift severity.

\section{Discussion}
The revised results change the interpretation of the study. The earlier unweighted ATCP approach was not reliable because the large source calibration set dominated the conformal threshold even though the source and target distributions were very different. SHIFT-ATCP corrected this problem by estimating source trust from shift metrics. Since both NHANES and Pima showed severe shift, the method assigned near-zero source weights and behaved close to target-only conformal. This is a useful automatic source-rejection behavior rather than a failure of the method.

The second major finding is that marginal coverage alone is not enough for clinical reliability. Safety-gated SHIFT-ATCP achieved near-nominal 90\% marginal coverage on both target datasets, but diabetes-positive coverage was much lower. DP-SHIFT-ATCP and risk-asymmetric SHIFT-ATCP improved diabetes-positive coverage substantially. This is clinically important because missing diabetes-positive patients may be more harmful than increasing referral for uncertain cases.

The trade-off is clear. SHIFT-ATCP is better for marginal coverage and efficiency. DP-SHIFT is better for diabetes-positive protection and class balance. Risk-asymmetric SHIFT provides the strongest diabetes-positive coverage but increases average set size and referral rate. Therefore, a clinically useful system should not report only ROC-AUC or marginal coverage. It should report marginal coverage, diabetes coverage, class coverage gap, referral rate, set size, and clinical cost together.

The source-trust stress test strengthens the novelty of the method. As simulated shift increased, domain AUC and PSI increased while source weight decreased. This supports the central assumption of SHIFT-ATCP: source calibration should be trusted only when source and target distributions are sufficiently close.

\subsection{Limitations}
The study used public datasets with different outcome definitions and measurement settings. NHANES diabetes status is survey-based and clinical-laboratory based, while Kaggle and Pima have their own dataset-specific definitions. Pima contains only adult female Pima Indian participants and lacks several variables available in the source dataset. The hyperparameter tuning was equal-budget and lightweight, not exhaustive. The clinical cost function was simple and should be adjusted with clinician input before deployment. Finally, DP-SHIFT improves diabetes-positive coverage by increasing uncertainty/referral, so it should be interpreted as a triage-support method rather than a replacement for diagnostic testing.

\section{Conclusion}
This work presents a revised and leakage-safe external transportability study for diabetes classification under dataset shift. Source-trained models were evaluated on NHANES and Pima using target calibration/test splits, split-specific ECE reporting, equal-budget tuning, shift-aware conformal baselines, and clinical reliability metrics. Conventional unweighted ATCP undercovered both targets. SHIFT-ATCP restored near-nominal marginal coverage by automatically down-weighting source calibration under severe shift. DP-SHIFT-ATCP improved diabetes-positive coverage, showing that class-wise clinical safety requires more than marginal conformal validity. The final recommendation is to use safety-gated SHIFT-ATCP when the main goal is marginal coverage and DP-SHIFT-ATCP when the clinical priority is protecting diabetes-positive patients.

\section*{Data Availability}
The source dataset is publicly available from Kaggle \cite{mustafa2023dataset}. The external validation dataset was constructed from publicly available NHANES 2017--2018 files \cite{nhanes2018overview}. The Pima Indians Diabetes dataset is publicly available \cite{smith1988using}.

\section*{Code Availability}
The reproducibility package refers to the project repository at \url{https://github.com/diwakardkk/Confirmal_Prediction}. The revised result folder used for this manuscript is \texttt{outputs\_dp\_shift\_atcp\_clinical\_novelty}.

\section*{Ethics Statement}
The analysis used public secondary datasets and did not involve direct participant contact or access to direct identifiers. Institution-specific and journal-specific requirements should be confirmed before submission.

\bibliographystyle{IEEEtran}
\bibliography{ref}

\begin{IEEEbiography}[{\includegraphics[width=1in,height=1.25in,clip,keepaspectratio]{a1.png}}]{Diwakar}
completed his Ph.D. in Computer Science from Babasaheb Bhimrao Ambedkar University, Lucknow (2026) and his M.Sc. in Computer Science from the University of Allahabad (2016). He works as an Assistant Professor at SRM University-AP, Andhra Pradesh. His areas of interest include artificial intelligence in healthcare applications, medical image processing, quantum computing and quantum machine learning, trustworthy and explainable AI, uncertainty-aware learning, and signal processing.
\end{IEEEbiography}

\begin{IEEEbiography}[{\includegraphics[width=1in,height=1.25in,clip,keepaspectratio]{a2.png}}]{Ansh Gupta}
is currently studying a bachelor's degree in Mathematical Sciences specializing in Data Analytics from Wenzhou-Kean University in Wenzhou, China, and he will be graduating in July 2027. His research areas include artificial intelligence in healthcare, medical image processing, cardiac MRI segmentation, ECG denoising, vision-language models, uncertainty quantification, and trustworthy AI.
\end{IEEEbiography}

\begin{IEEEbiography}[{\includegraphics[width=1in,height=1.25in,clip,keepaspectratio]{a3.jpeg}}]{Gaurav Gupta}
received the Ph.D. degree in Mathematics from Hemwati Nandan Bahuguna Garhwal University, Srinagar, India, and the M.Sc. degree in Mathematics from Gurukul Kangri Vishwavidyalaya, Haridwar, India. He is currently an Associate Professor of Mathematical Sciences with the College of Science, Mathematics and Technology at Wenzhou-Kean University, Wenzhou, China.
\end{IEEEbiography}

\EOD
\end{document}
